% TOP_PROBLEM: {"id": 30002739, "title": "Periodic Gibbs Phases in Two Dimensions", "category_id": 16, "category": {"id": 16, "name": "physics", "display_name": "Mathematical Physics", "description": "Problems at the intersection of mathematics and physics.", "slug": "physics", "order_index": 16, "created_at": "2026-07-31T15:26:25.671Z"}, "status": "open", "statusQualification": "The allowed state space and interaction assumptions must be those of the report. In particular, the source statement should not be generalized without review to all continuous-spin, noncompact, or arbitrary inhomogeneous interactions."}
% FIRST_POSTED: 2026-10-11
% ENTRY_KIND: statement_only
% SUBMISSION: {"kind": "statement_only", "category": "top_problems", "problem_id": "30002739", "proposal_type": "established", "source": "Mini-Workshop: Dynamical versus Diffraction Spectra in the Theory of Quasicrystals (2014), Oberwolfach Reports, DOI 10.4171/owr/2014/53, pp. 2987--3013. Primary report accessed 2026-10-11; relevant PDF pages 9, 10. https://doi.org/10.4171/owr/2014/53; Koch and Radin, positive-temperature quasicrystal model (catalogue-listed context; not independently consulted in this preparation). https://doi.org/10.1007/s10955-009-9896-9", "status_evidence": "Canonical source status open retained. Supplied literature triage dated 2026-08-21: No general theorem was found proving that every two-dimensional finite-range positive-temperature extremal Gibbs state is periodic; known results and models leave the proposed no-quasicrystal statement unresolved. Primary-report access checked 2026-10-11; this historical report and inherited triage do not certify current resolution of every component.", "status_checked_on": "2026-10-11", "duplicate_check": "Compared canonical ID, title and normalized target with the complete local UnsolvedMath registry; normalized exact-text matches: none; title matches: none. Related or overlapping canonical records: none identified in this batch. Primary source context and the existing OWR statement audit used for scope and duplicate review; an exhaustive mathematical-equivalence/literature search was not performed.", "attribution": "Original proposer not named in the supplied canonical record; attribution remains to the identified Oberwolfach contribution/report, not its editors as invented proposers. Explanatory statement preparation and source/semantic audit assisted by OpenAI GPT-6 Codex; no human mathematical authorship or independent proof verification is claimed."}
% SOURCE_STATEMENT: "At positive temperature, is every extremal Gibbs measure for a finite-range interaction in two dimensions periodic?"
% Preparation and final self-audit: CONTRIBUTING.md prompts C and D; source/audit artifacts retained in the local physics_statements_2026_10_11 batch.
% Imported from: physics_statements.tex; UnsolvedMath ID 30002739
% !TeX program = lualatex
\documentclass[11pt]{article}
\usepackage[margin=25mm]{geometry}
\usepackage{fontspec}
\setmainfont{Latin Modern Roman}
\usepackage{amsmath,amssymb,mathtools}
\usepackage{unicode-math}
\setmathfont{Latin Modern Math}
\usepackage{xurl,hyperref}
\hypersetup{colorlinks=true,urlcolor=blue}
\setcounter{secnumdepth}{0}
\setlength{\parindent}{0pt}
\setlength{\parskip}{5pt}
\setlength{\emergencystretch}{3em}
\newcommand{\sref}[2]{\hyperlink{source:#1:#2}{[S#2]}}
\newcommand{\cataloguescope}[1]{\paragraph{Recorded target.}#1}
\newcommand{\R}{\mathbb R}
\newcommand{\C}{\mathbb C}
\newcommand{\N}{\mathbb N}
\newcommand{\Z}{\mathbb Z}
\newcommand{\Q}{\mathbb Q}
\newcommand{\F}{\mathbb F}
\newcommand{\eps}{\varepsilon}
\begin{document}
% UnsolvedMath ID 30002739; source catalogue code OWR-13355-003
\setcounter{section}{30002738}
\section{Periodic Gibbs Phases in Two Dimensions}\label{30002739}
{\small UnsolvedMath ID 30002739 · Mathematical Physics}
\subsection{Definitions and mathematical statement}

\paragraph{Shared notation.}$\mathbb N=\{1,2,\ldots\}$; the standard number fields have their usual meanings. All additional conventions are stated locally.


\paragraph{Definitions and conventions.}
For a lattice interaction of finite range, a Gibbs measure satisfies the Dobrushin--Lanford--Ruelle conditional-probability equations at inverse temperature $\beta\in(0,\infty)$. An extremal Gibbs measure is not a nontrivial convex mixture of other Gibbs measures for that same interaction. Periodicity means invariance under a finite-index subgroup of lattice translations, allowing a finite unit cell. Periodic is stronger than merely having a periodic-looking local correlation pattern.

\paragraph{Target retained from the canonical source catalogue.}
\paragraph{Catalogue wording retained for traceability.}
At positive temperature, is every extremal Gibbs measure for a finite-range interaction in two dimensions periodic?

\paragraph{Source scope and qualification.}
The allowed state space and interaction assumptions must be those of the report. In particular, the source statement should not be generalized without review to all continuous-spin, noncompact, or arbitrary inhomogeneous interactions.
\subsection{Short English statement}
Must every pure equilibrium phase of a finite-range two-dimensional system at positive temperature be periodic?
\subsection{Sources}
\begin{itemize}
\item[\hypertarget{source:30002739:1}{S1}] ULAM.AI and the UnsolvedMath Contributors, UnsolvedMath, record 30002739 (OWR-13355-003). Catalogue curation under CC BY 4.0; definitions and English exposition added locally.
\url{https://huggingface.co/datasets/ulamai/UnsolvedMath}.
\item[\hypertarget{source:30002739:2}{S2}] Mini-Workshop: Dynamical versus Diffraction Spectra in the Theory of Quasicrystals (2014), Oberwolfach Reports, DOI 10.4171/owr/2014/53, pp. 2987--3013. Primary report accessed 2026-10-11; relevant PDF pages 9, 10.
\url{https://doi.org/10.4171/owr/2014/53}.
\item[\hypertarget{source:30002739:3}{S3}] Koch and Radin, positive-temperature quasicrystal model (catalogue-listed context; not independently consulted in this preparation).
\url{https://doi.org/10.1007/s10955-009-9896-9}.
\end{itemize}
\label{end:30002739}
\end{document}
